\part{Securing Devices}

Unit 3 secured the network the devices talk over. Unit 4 secures the devices themselves, on the
premise that a perfectly locked-down network still fails if the laptop connected to it is running
outdated software with a known, unpatched vulnerability. Chapter~\ref{ch:4.1} surveys where those
vulnerabilities actually come from -- unpatched operating systems, insecure default
configurations, outdated firmware -- and why an endpoint is often the weakest link precisely
because it is the device a user controls directly.

Chapter~\ref{ch:4.2} takes on the first line of defense for any device: authentication. This
chapter goes past ``use a strong password'' into what actually makes multi-factor authentication
resistant to the credential-theft techniques covered back in Unit 1, and where it can still be
defeated. Chapter~\ref{ch:4.3} broadens out to endpoint hardening generally -- disabling unused
services, enforcing least privilege on the device itself, keeping software patched -- the
device-level version of the access-control philosophy introduced in Chapter~\ref{ch:2.3}.

Chapter~\ref{ch:4.4} closes the unit, again, with auditing: how an organization confirms its
endpoints are actually configured the way policy says they should be, rather than assuming
compliance and finding out otherwise during an incident. By now the prevent-control-detect
rhythm from the last two units should feel automatic, which is the point -- it is the same
framework, one layer deeper.
